% Equivalent reduction promoted after complete root and reduction_descent audits. Original candidate retained unchanged.
\subsection{A determinant-variation expression for the remaining rank map}

\begin{lemma}[The actual symplectic cohomology sequence]
\label{lem:DE-A3-positive-monad-note}
\status{Proved equivalent reduction; complete independent root and reduction/descent proof audits passed.}
Keep the notation of the acyclic-factorization lemma and put
$\mathcal M=\mathcal O_T^2\oplus\mathfrak g^\vee$. There is an exact
sequence of bundles
\[
0\longrightarrow Q^\vee\xrightarrow{i=(a^\vee,f^\vee)}
\mathcal M\xrightarrow{p=(aJ_2,fJ_{\mathfrak g})}Q
\longrightarrow0.                                      \tag{DE-MN1}
\]
The composite is zero by (DE-TZ1). Its middle bundle is fixed as
the sections $s_0,s_1$ vary. Its rank-two image is Lagrangian for
the alternating pairing $\mathcal M\to\mathcal M^\vee$
with matrix $\operatorname{diag}(J_2,J_{\mathfrak g})$.
At a sharp source it gives
\[
0\longrightarrow H^1(Q^\vee)\longrightarrow
H^1(\mathcal O_T)^2\longrightarrow H^1(Q)\longrightarrow0.
                                                               \tag{DE-MN2}
\]
Also $\operatorname{coker}(f)=i_*A$, canonically up to the fixed
determinant identifications.
\end{lemma}
\begin{proof}
In the earlier chart, $a$ has columns $(1,0),(uh^2,w)$ and $f$
has columns $(uw,0),(0,1)$. Their images together generate $Q$
even on $D$. Since $J_2,J_{\mathfrak g}$ are invertible, $p$ is
surjective. Its dual, composed with the perfect alternating form on
$\mathcal M$, identifies its rank-two kernel with the displayed
image of $i$; the composite is already zero, and this identification
can equally be checked in the same chart. This proves exactness
and the Lagrangian assertion.

The cokernel of $f$ is locally the first-coordinate line on $D$.
The two sections induced by $a$ are the row $(1,uh^2)$, precisely
the squared pencil. This identifies that line with $A$ on overlaps.
As $\mathfrak g$ is acyclic, $H^j(Q)=H^j(D,A)$.
At a sharp source these spaces have dimension two for $j=0,1$
and vanish for $j=2$. The map $H^0(\mathcal M)=k^2\to H^0(Q)$
is an isomorphism: multiplication by $J_2$ only interchanges the
two known basis sections. Its Serre dual makes
$H^2(Q^\vee)\to H^2(\mathcal M)$ an isomorphism.
The long cohomology sequence then gives (DE-MN2).
\end{proof}

\begin{lemma}[Determinant variation gives the cohomological tangent map]
\label{lem:DE-A3-positive-monad-variation-note}
\status{Proved equivalent reduction; complete independent root and reduction/descent proof audits passed.}
On the sharp disjoint-zero-scheme parameter open, let
$S_s\subset H^1(\mathcal O_T)^2$ be the Lagrangian subspace in
(DE-MN2). For a section direction $t=(t_0,t_1)$ put
\[
 \dot w=s_0\wedge t_1+t_0\wedge s_1\in H^0(T,LP).
\]
Under (DE-MN2), the derivative of the subspace $S_s$ is the map
\[
 H^1(Q^\vee)\xrightarrow{H^1(\dot w J_Q)}H^1(Q),        \tag{DE-MN3}
\]
where $J_Q:Q^\vee\to Q\otimes(\det Q)^{-1}$ is the canonical
inverse alternating identification. All bundle determinants here
use $\det Q=LP$.
Equivalently, after the already specified graph identifications
using the torsor classes $\alpha,\beta$, this is the derivative
of the actual cup matrix $T_s$.
\end{lemma}
\begin{proof}
For an infinitesimal deformation of a Lagrangian subbundle $i_s$
of the fixed bundle $\mathcal M$, its derivative is represented
by $p\dot i:Q^\vee\to Q$. This map is independent of a local
choice of frame for the deformed $Q^\vee$: changing such a frame
adds a term $iX$ to $\dot i$, killed by $p$.

Choose a local determinant-one frame of $V=E_0$ normalizing the
nonvanishing section $s_0$ to $(1,0)$, also over the dual numbers.
The induced first-order change of frame on
$\operatorname{End}^0(V)/\mathcal O_T$ is zero. Indeed the
infinitesimal change is the commutator $[X,Y]$ of two trace-zero
two-by-two matrices, and in characteristic two this commutator
is scalar. Equivalently the rank-two quotient representation
of $\operatorname{SL}_2$ is its Frobenius twist. The frame of $P$
is fixed, so the frame of $\mathfrak g$ is stationary to first
order. The chosen determinant-one frame also leaves the local
coordinate of the section $w$ of $LP$ unchanged.

In this frame the two matrices retain their earlier form.
Because $d(h^2)=0$ and the local unit $u$ is parameter-independent,
\[
 \dot a=\begin{pmatrix}0&0\\0&\dot w\end{pmatrix},
 \qquad
 \dot f=\begin{pmatrix}u\dot w&0\\0&0\end{pmatrix}.
\]
Consequently
\[
 p\dot i=aJ_2\dot a^t+f(u^{-1}J_2)\dot f^t
        =\begin{pmatrix}0&\dot w\\\dot w&0\end{pmatrix}.
\]
The expression is intrinsic, so the same computation on the other
nonvanishing chart proves globally that $p\dot i=\dot wJ_Q$.

Finally differentiate a Cech representative of a section of the
cohomology subspace $S_s$. Terms caused by differentiating its
representative inside $H^1(Q^\vee)$ still lie in $S_s$ and
disappear in the quotient. The remaining term is the cohomology
map induced by $p\dot i$. The constant dimensions on the sharp
open justify cohomology base change; (DE-MN2) identifies the
quotient with $H^1(Q)$. This proves (DE-MN3).
In graph coordinates $S_s=\{\alpha v+\beta T_s v\}$, its
derivative is $v\mapsto\beta(dT_s)v$ modulo $S_s$, giving the
last statement with the same basis interchange as in (DE-TZ3).
\end{proof}

\begin{remark}[Precise unproved lower bound]
\label{rem:DE-A3-positive-monad-gate-note}
\status{Proved scope statement; the required variation-image lower bound remains unproved.}
Put
\[
\begin{split}
\Delta_s&=\operatorname{im}\bigl[H^0(E_0)\oplus H^0(E_1)
\xrightarrow{(t_0,t_1)\mapsto s_0\wedge t_1+t_0\wedge s_1}
H^0(LP)\bigr],\\
\Psi_s:\Delta_s&\longrightarrow
\operatorname{Hom}_k\bigl(H^1(Q^\vee),H^1(Q)\bigr),
\qquad \dot w\longmapsto H^1(\dot wJ_Q).
\end{split}
\]
The uniform sharpness problem would follow from
$\dim_k\operatorname{im}\Psi_s\ge2$ at one suitable sharp source.
This is a condition on the span of the entire family of variation
maps, not a rank assertion about one fixed member of that family.
The acyclicity of $\mathfrak g$,
the exactness of (DE-MN1), and the trace-zero theorem do not by
themselves imply that lower bound. No Petri multiplication
identification or generic dimension claim is made in this note.
\end{remark}
